% ============================================================
\chapter{G\'en\'eration de Texte}
\label{ch:generation}
% ============================================================

Un mod\`ele de langue qui pr\'edit le mot suivant est, litt\'eralement, un g\'en\'erateur de texte en puissance. Mais entre la distribution $P(x_t \mid x_{<t})$ et un texte coh\'erent, fluide et informatif, il y a un gouffre --- celui de la \emph{strat\'egie de d\'ecodage}. Prendre syst\'ematiquement le mot le plus probable (d\'ecodage glouton) produit un texte r\'ep\'etitif et ennuyeux. \'Echantillonner na\"ivement dans la distribution produit un texte incoh\'erent. Les solutions --- beam search, top-$k$, nucleus sampling (top-$p$), temp\'erature --- sont autant de compromis entre \emph{qualit\'e} et \emph{diversit\'e}. Ari Holtzman et ses collaborateurs, dans leur article ``The Curious Case of Neural Text Degeneration'' (2020), ont montr\'e que le nucleus sampling r\'esout \'el\'egamment le probl\`eme de d\'eg\'en\'erescence. Ce chapitre explore ces strat\'egies et leurs fondements probabilistes.

\begin{intuition}[title=\textbf{Du mod\`ele de langue au g\'en\'erateur}]
Un mod\`ele de langue auto-r\'egressif d\'efinit $P(x_t \mid x_{<t})$.
G\'en\'erer du texte revient \`a \'echantillonner s\'equentiellement dans cette
distribution. Le choix de la strat\'egie de d\'ecodage influence radicalement la
qualit\'e, la diversit\'e et la coh\'erence du texte produit.
\end{intuition}

% ------------------------------------------------------------
\section{D\'ecodage glouton}
\label{sec:gen:greedy}
% ------------------------------------------------------------

\begin{definition}[D\'ecodage glouton]
Le d\'ecodage glouton (\emph{greedy decoding}) s\'electionne \`a chaque pas
le token le plus probable :
\[
x_t = \argmax_{w \in \mathcal{V}} P(w \mid x_{<t})
\]
\end{definition}

\begin{attention}[title=\textbf{Limites du d\'ecodage glouton}]
Le d\'ecodage glouton ne garantit pas la s\'equence globalement optimale. Il
peut produire des r\'ep\'etitions et manque de diversit\'e. La s\'equence la
plus probable n'est pas toujours la plus naturelle.
\end{attention}

% ------------------------------------------------------------
\section{Recherche en faisceau (Beam Search)}
\label{sec:gen:beam}
% ------------------------------------------------------------

\begin{definition}[Beam search]
La recherche en faisceau maintient les $B$ s\'equences partielles les plus
probables (\emph{beams}) \`a chaque pas de d\'ecodage :
\[
\mathcal{B}_t = \text{top-}B\left\{(s \oplus w, \text{score}(s) + \log P(w \mid s)) : s \in \mathcal{B}_{t-1}, w \in \mathcal{V}\right\}
\]
o\`u $B$ est la largeur du faisceau.
\end{definition}

\begin{proposition}[Propri\'et\'es du beam search]
\begin{enumerate}
    \item Pour $B = 1$ : \'equivalent au d\'ecodage glouton.
    \item Pour $B = \abs{\mathcal{V}}^T$ : \'equivalent \`a la recherche exhaustive.
    \item En pratique, $B \in [4, 10]$ offre un bon compromis.
\end{enumerate}
\end{proposition}

\subsection{Normalisation par la longueur}

Le beam search favorise les s\'equences courtes (plus de termes $\log P < 0$
pour les longues). On normalise le score par la longueur :
\[
\text{score}_{\text{norm}}(s) = \frac{1}{\abs{s}^\alpha} \sum_{t=1}^{\abs{s}} \log P(x_t \mid x_{<t})
\]
o\`u $\alpha \in [0.6, 1.0]$ est un hyperparam\`etre (Wu et al., 2016).

% ------------------------------------------------------------
\section{\'Echantillonnage stochastique}
\label{sec:gen:sampling}
% ------------------------------------------------------------

\begin{definition}[\'Echantillonnage avec temp\'erature]
L'\'echantillonnage avec temp\'erature $\tau > 0$ modifie la distribution :
\[
P_\tau(w \mid x_{<t}) = \frac{\exp(\text{logit}_w / \tau)}{\sum_{w'} \exp(\text{logit}_{w'} / \tau)}
\]
\begin{itemize}
    \item $\tau \to 0$ : converge vers le d\'ecodage glouton.
    \item $\tau = 1$ : distribution originale.
    \item $\tau > 1$ : distribution plus uniforme (plus de diversit\'e).
\end{itemize}
\end{definition}

\begin{definition}[Top-$k$ sampling]
Le top-$k$ sampling (Fan et al., 2018) restreint l'\'echantillonnage aux $k$
tokens les plus probables :
\[
P_{\text{top-}k}(w \mid x_{<t}) = \begin{cases}
\frac{P(w \mid x_{<t})}{\sum_{w' \in \mathcal{V}_k} P(w' \mid x_{<t})} & \text{si } w \in \mathcal{V}_k \\
0 & \text{sinon}
\end{cases}
\]
o\`u $\mathcal{V}_k$ contient les $k$ tokens les plus probables.
\end{definition}

\begin{definition}[Top-$p$ sampling (nucleus)]
Le top-$p$ sampling (Holtzman et al., 2020) s\'electionne le plus petit
ensemble $\mathcal{V}_p$ tel que :
\[
\sum_{w \in \mathcal{V}_p} P(w \mid x_{<t}) \geq p
\]
puis \'echantillonne uniform\'ement dans $\mathcal{V}_p$ (apr\`es
renormalisation).
\end{definition}

\begin{bonnepratique}[title=\textbf{Combinaison des strat\'egies}]
En pratique, on combine souvent temp\'erature, top-$k$ et top-$p$. Par exemple,
$\tau = 0.7$, $k = 50$, $p = 0.9$ offre un bon \'equilibre entre coh\'erence
et diversit\'e pour la g\'en\'eration cr\'eative.
\end{bonnepratique}

% ------------------------------------------------------------
\section{P\'enalit\'es de r\'ep\'etition}
\label{sec:gen:repetition}
% ------------------------------------------------------------

\begin{definition}[P\'enalit\'e de r\'ep\'etition]
Pour \'eviter les r\'ep\'etitions, on p\'enalise les tokens d\'ej\`a g\'en\'er\'es :
\[
\text{logit}'_w = \begin{cases}
\text{logit}_w / \theta & \text{si } w \in x_{<t} \text{ et } \text{logit}_w > 0 \\
\text{logit}_w \cdot \theta & \text{si } w \in x_{<t} \text{ et } \text{logit}_w \leq 0
\end{cases}
\]
o\`u $\theta > 1$ est le facteur de p\'enalit\'e.
\end{definition}

% ------------------------------------------------------------
\section{D\'ecodage contrast\'e}
\label{sec:gen:contrastive}
% ------------------------------------------------------------

\begin{definition}[Contrastive decoding]
Le d\'ecodage contrast\'e (Li et al., 2023) s\'electionne les tokens qui
maximisent la diff\'erence de log-probabilit\'e entre un mod\`ele expert
$\pi_{\text{exp}}$ et un mod\`ele amateur $\pi_{\text{ama}}$ :
\[
x_t = \argmax_{w \in \mathcal{V}_p} \left[\log \pi_{\text{exp}}(w \mid x_{<t}) - \alpha \log \pi_{\text{ama}}(w \mid x_{<t})\right]
\]
\end{definition}

% ------------------------------------------------------------
\section{D\'ecodage sp\'eculatif}
\label{sec:gen:speculative}
% ------------------------------------------------------------

\begin{definition}[Speculative decoding]
Le d\'ecodage sp\'eculatif (Leviathan et al., 2023) utilise un petit mod\`ele
rapide (\emph{draft model}) pour proposer $K$ tokens, puis un grand mod\`ele
les v\'erifie en parall\`ele. Les tokens accept\'es sont ceux dont la probabilit\'e
sous le grand mod\`ele est suffisante, ce qui acc\'el\`ere la g\'en\'eration
sans changer la distribution.
\end{definition}

% ------------------------------------------------------------
\section{M\'etriques de qualit\'e de g\'en\'eration}
\label{sec:gen:metriques}
% ------------------------------------------------------------

\begin{formules}[title=\textbf{M\'etriques de g\'en\'eration}]
\begin{align}
\text{BLEU} &= \text{BP} \cdot \exp\!\left(\sum_{n=1}^{N} w_n \log p_n\right) & \text{(pr\'ecision $n$-grammes)} \\
\text{ROUGE-L} &= F_1(\text{LCS}) & \text{(plus longue sous-s\'equence commune)} \\
\text{METEOR} &= F_1 \cdot (1 - \text{Penalty}) & \text{(alignement flexible)} \\
\text{BERTScore} &= F_1(\text{sim cosinus BERT}) & \text{(similarit\'e contextuelle)}
\end{align}
\end{formules}

% ------------------------------------------------------------
\section{Impl\'ementation}
\label{sec:gen:implementation}
% ------------------------------------------------------------

\begin{codeblock}[title=\textbf{Strat\'egies de d\'ecodage avec Hugging Face}]
\begin{minted}{python}
from transformers import AutoTokenizer, AutoModelForCausalLM

model_name = "gpt2-medium"
tokenizer = AutoTokenizer.from_pretrained(model_name)
model = AutoModelForCausalLM.from_pretrained(model_name)

prompt = "The future of artificial intelligence"
input_ids = tokenizer.encode(prompt, return_tensors="pt")

# Greedy decoding
greedy = model.generate(input_ids, max_length=50)

# Beam search
beam = model.generate(input_ids, max_length=50, num_beams=5,
                      length_penalty=0.8, no_repeat_ngram_size=3)

# Top-k + Top-p sampling
sampled = model.generate(input_ids, max_length=50, do_sample=True,
                         temperature=0.7, top_k=50, top_p=0.9,
                         repetition_penalty=1.2)

for name, ids in [("Greedy", greedy), ("Beam", beam), ("Sampled", sampled)]:
    print(f"{name}: {tokenizer.decode(ids[0], skip_special_tokens=True)}")
\end{minted}
\end{codeblock}

\begin{codeblock}[title=\textbf{D\'ecodage pas \`a pas}]
\begin{minted}{python}
import torch
import torch.nn.functional as F

def generate_top_p(model, input_ids, max_new_tokens=50, temperature=0.8,
                   top_p=0.9):
    generated = input_ids.clone()
    for _ in range(max_new_tokens):
        with torch.no_grad():
            outputs = model(generated)
            logits = outputs.logits[:, -1, :] / temperature

        # Top-p filtering
        sorted_logits, sorted_indices = torch.sort(logits, descending=True)
        cumulative_probs = torch.cumsum(F.softmax(sorted_logits, dim=-1), dim=-1)
        mask = cumulative_probs - F.softmax(sorted_logits, dim=-1) >= top_p
        sorted_logits[mask] = float('-inf')
        logits.scatter_(1, sorted_indices, sorted_logits)

        probs = F.softmax(logits, dim=-1)
        next_token = torch.multinomial(probs, num_samples=1)
        generated = torch.cat([generated, next_token], dim=-1)

        if next_token.item() == tokenizer.eos_token_id:
            break
    return generated
\end{minted}
\end{codeblock}

% ------------------------------------------------------------
\section{G\'en\'eration contr\^ol\'ee}
\label{sec:gen:controlee}
% ------------------------------------------------------------

\begin{definition}[G\'en\'eration contr\^ol\'ee]
La \emph{g\'en\'eration contr\^ol\'ee} vise \`a orienter la sortie d'un mod\`ele
de langue selon des attributs sp\'ecifi\'es (ton, sujet, style) sans
r\'e-entra\^inement. Approches :
\begin{itemize}
    \item \textbf{PPLM} (Dathathri et al., 2020) : perturbation des gradients.
    \item \textbf{Classifier-free guidance} : interpolation entre distributions
          conditionn\'ee et non conditionn\'ee.
    \item \textbf{Constrained decoding} : contraintes lexicales dures.
\end{itemize}
\end{definition}

% ------------------------------------------------------------
\section{Exercices}
\label{sec:gen:exercices}
% ------------------------------------------------------------

\begin{exercice}[$\star$]
Montrez que le d\'ecodage glouton est un cas particulier du beam search avec
$B = 1$.
\end{exercice}

\begin{exercice}[$\star$]
Calculez la distribution top-$p$ pour $p = 0.9$ \'etant donn\'e les
probabilit\'es $(0.4, 0.3, 0.15, 0.1, 0.05)$.
\end{exercice}

\begin{exercice}[$\star\star$]
Impl\'ementez le beam search avec normalisation par la longueur. Testez avec
diff\'erentes valeurs de $\alpha$ et $B$ sur GPT-2.
\end{exercice}

\begin{exercice}[$\star\star$]
Comparez quantitativement (BLEU, diversit\'e) la g\'en\'eration par greedy,
beam search ($B=5$), et nucleus sampling ($p=0.9$) sur un corpus de test.
\end{exercice}

\begin{exercice}[$\star\star\star$]
Impl\'ementez le d\'ecodage sp\'eculatif avec un mod\`ele draft (GPT-2 small)
et un mod\`ele cible (GPT-2 large). Mesurez l'acc\'el\'eration obtenue et
v\'erifiez que la distribution de sortie est pr\'eserv\'ee.
\end{exercice}
